\documentclass[12pt]{article}
\usepackage[utf8]{inputenc}
\usepackage[margin=0.75in]{geometry}
\usepackage{amsmath}
\usepackage{amssymb}
\usepackage{amsthm}
\usepackage{graphicx} % Required for inserting images
\usepackage{array}

\title{MAT1100}
\author{Aaron Avram}
\date{September 8th 2026}

\begin{document}
\theoremstyle{plain}
\newtheorem{theorem}{Theorem}
\newtheorem{proposition}{Proposition}
\newtheorem{lemma}[theorem]{Lemma}
\newtheorem{corollary}[theorem]{Corollary}
\newtheorem{claim}[theorem]{Claim}

\theoremstyle{definition}
\newtheorem{definition}{Definition}
\newtheorem{example}{Example}
\newtheorem{exercise}[theorem]{Exercise}

\theoremstyle{remark}
\newtheorem{remark}[theorem]{Remark}
\newtheorem{note}[theorem]{Note}

\maketitle

\tableofcontents
\newpage

\section{Group Theory}

\subsection{Basic Definitions and Results}
\begin{definition}
    A Group is a set with a binary operation $\star: G \times G \to G$ satisfying
    \begin{enumerate}
        \item Associativity: for all $a, b, c \in G$, $(a \star b) \star c = a \star (b \star c)$
        \item Existence of an identity: $\exists e \in G$ such that for all $a \in G$, $a = a \star e = e \star a$.
        \item Existence of inverses: for all $a \in G$ there exists $b \in G$ such that $ab = ba = e$.
    \end{enumerate}
\end{definition}

\begin{example}
    Some examples
    \begin{enumerate}
        \item $(\mathbb Z, +)$
        \item $GL_n(\mathbb K) = \{ A \in M_n(\mathbb K) | \det(A) \neq 0 \}$
        \item $SL_n(\mathbb K) = \{ A \in M_n(\mathbb K) | \det(A) = 1 \}$
    \end{enumerate}
\end{example}
\begin{example}
    $D_8$ the group of symmetries of the unit square S.
    \[ D_8 = \{ T \in GL_2(\mathbb R) : T(S) = S \wedge \| T(x) - T(y) \| = \| x - y \|, \forall x, y \in \mathbb R^2 \} \]
    it is not hard to check that this forms a group under function composition.
\end{example}
\begin{definition}
    Subgroups $H \subset G$ is a subgroup and denoted $H \leq G$ if it is a group
    under the operation of $G$.

    Equivalently:
    \begin{enumerate}
        \item $e \in H$
        \item $xy \in H$ for all $x, y \in H$
        \item $x^{-1} \in H$ for all $x \in H$.
    \end{enumerate}
\end{definition}
\begin{remark}
    There is a subgroup criterion that is sometimes useful:
    $H \subseteq G$ is a subgroup if and only if $H \neq \varnothing$
    and for all $x, y \in H$, $xy^{-1} \in H$
\end{remark}
\begin{example}
    Some examples
    \begin{enumerate}
        \item $SL_n(\mathbb K) \leq GL_n(\mathbb K)$
        \item $(\mathbb Z, +) \leq (\mathbb Q, +)$
        \item $D_8 \leq GL_2(\mathbb R)$.
    \end{enumerate}
\end{example}

\begin{definition}
    For $g \in G$ a group, the cyclic subgroup generated by $g$ is
    \[ \langle g \rangle = \{ g^n \mid n \in \mathbb Z \} \]
\end{definition}

\begin{example}
    $S_3 := \{ \sigma: \{ 1, 2, 3 \} \to \{ 1, 2, 3 \} \text{ a bijection} \}$
    \[ \langle (1, 2) \rangle = \{ e, (1, 2) \} \]
    and
    \[ \langle (1, 2, 3) \rangle = \{ e, (1, 2, 3), (1, 3, 2) \} \]
\end{example}

\begin{definition}
    The order of a group $G$ is its cardinality and denoted $|G|$.
    For $g \in G$, $|g| = \min \{ n \in \mathbb N : g^n = e \}$
    called the order of $g$.
\end{definition}

\begin{proposition}
    For $g \in G$ we have $|\langle g \rangle| = |g|$.
\end{proposition}
\begin{proof}
    Set $n = |g|$. By the minimality of the order $e, g, \ldots, g^{n-1}$ are all distinct.
    Now it suffices to show that $\langle g \rangle  = \{ e, g, \ldots, g^{n-1} \}$
    the reverse containment is clear. For the forward containment, let $g^a \in \langle g \rangle $.
    $a \in \mathbb Z$, so there exists $k \in \mathbb Z$ and $r \in {0, \ldots, n-1}$ such that
    \[ a = kn + r \]
    then
    \[ g^a = g^{kn + r} = g^{kn}g^r = g^r \]
    hence $g^a$ is represented by an element in the right hand set which proves the proposition.
\end{proof}

\begin{definition}
    A group homomorphism is a map $\varphi: G \to G'$ such that
    $\varphi(gh) = \varphi(g)\varphi(h)$ for all $g, h \in G$.

    \begin{enumerate}
    \item The kernel of $\varphi$: $\ker \varphi = \{ g \in G | \varphi(g) = e_{G'} \}$
    \item The image of $\varphi$: $\operatorname{Im} \varphi = \{ \varphi(g) \mid g \in G \}$
    \end{enumerate}
\end{definition}

\begin{proposition}
    $\ker \varphi \leq G$ and $\operatorname{Im} \varphi \leq H$.
\end{proposition}

\begin{example}
        $\varphi = \det: GL_n(\mathbb K) \to \mathbb K^\times$ with $\ker \varphi = SL_n(\mathbb K)$.
\end{example}

\begin{example}
    Now we compute $D_8$ explicitly. Label the vertices of the square $1, 2, 3, 4$ starting from the bottom left and moving counter clockwise.

    A symmetry of the square must be a permutation as the vertices are exactly the points of the square
    which have another point a distance of the diagonal away and isometries must preserve this.

    Additionally isometries must preserve adjacency

    Hence given an element $T \in D_8$, $T(1)$ has 4 possibilities,
    then given $T(1)$, $T(2)$ has 2 and then by adjacency the isometry is now determined.

    Therefore $|D_8| \leq 8$.

    Now it is clear that two isometries are:
    \begin{enumerate}
        \item s: reflection about the diagonal 1-3.
        \item r: counter-clockwise rotation by $\frac{\pi}{2}$
    \end{enumerate}
    and from here it is easy to show that
    \[ D_8 = \{ e, s, r, r^2, r^3, rs, r^2s, r^3s \} \]
    and from here we recover another group as a subgroup
    \[ \{ e, s, r^2, r^2s \} = K_4 \]
    the Klein 4-group.
\end{example}

\subsection{Group Actions}
\begin{definition}
    The actiion of a group $G$ on a set $X \neq \varnothing$ is a map
    \[ G \times X \to X \]
    given by
    \[ (g, x) \mapsto g \cdot x \]
    such that 
    \[ e \cdot x = x \quad \forall x \in X \]
    and
    \[ (gh) \cdot x = g \cdot (h \cdot x) \forall x \in X \]
\end{definition}
\begin{example}
    \begin{enumerate}
        \item[]
        \item[(1)]
        $D_8 \curvearrowright V$ where $V = \{ 1, 2, 3, 4 \}$ are the vertices of the unit square.
        This action is given by:
        \[ (g, i) \mapsto g \cdot i = \sigma_g(i) \]
        where $\sigma_g$ is the permutation of the vertices associated with $g$.

        \item[(2)]
        $D_8 \curvearrowright \Delta = \{ d_1, d_2 \}$
        \[ g \cdot \{i, j\} = \{ g(i), g(j) \} \]
        by adjacency.
    \end{enumerate}
\end{example}
\begin{theorem}
    Let $G \curvearrowright X$ be an action of a group $G$ on a set $X$.
    Then for each $g \in G$ the map
    \[ \rho(g): X \to X \]
    given by
    \[ \rho(g)(x) = g \cdot x \]
    is a bijection with inverse $\rho(g^{-1})$ and the map
    \[ \rho: G \to \operatorname{Sym}(X) \]
    given by:
    \[ \rho: g \to \rho(g) \]
    is a group homomorphism.
\end{theorem}
\begin{remark}
    The converse is also true, a homomorphism $\rho: G \to \operatorname{Sym}(X)$
    induces an action.
\end{remark}
\begin{definition}
    For an action $G \to X$ the orbit of $x \in X$ is $\mathcal O_x = \{ g \cdot x : g \in G \}$.
    An action is transitive if it has only one orbit.
\end{definition}
\begin{proposition}
    The relation on $X$ given by $a \sim b \iff \exists g \in G$ s.t. $g \cdot a = b$
    is an equivalence relation, and its equivalence classes are the distinct orbits of $X$.
\end{proposition}
\begin{proposition}
    A subgroup $H \leq G$ induces an action $H \curvearrowright G$ given by
    \[ h \cdot g = gh^{-1} \]
    and its orbits
    \[ \mathcal O_g = \{ h \cdot g : h \in H \} = \{ gh^{-1} : h \in H \} = gH \]
    is called the left coset of $g$ in $H$.
    It is clear $|gH| = |H|$.
\end{proposition}
\begin{definition}
    Write $G/H$ to be the set of left cosets, then
    \[ [G: H] = \#(G/H) \]
\end{definition}
\begin{theorem}
    If $H \leq G$ we have $|G| = |H| [G:H]$. In particular the order of a subgroup
    divides the order of $G$
\end{theorem}
\begin{corollary}
    The order of an element divides the order of the group.
\end{corollary}
\end{document}
